\documentclass[11pt]{article}

\usepackage[margin=1in]{geometry}
\usepackage{times}
\usepackage{hyperref}
\usepackage{url}
\usepackage{enumitem}
\usepackage{titlesec}

% Compact spacing
\setlength{\parskip}{2pt}
\setlength{\parindent}{0pt}
\titlespacing*{\section}{0pt}{6pt}{3pt}
\titlespacing*{\subsection}{0pt}{4pt}{2pt}

\title{\vspace{-0.6in}\Large Do Robots Need Cognitive Memory? An Empirical Study of\\Multi-Tier LLM-Augmented Memory for Robotic Agents\vspace{-0.15in}}

\author{Alex Schott \and Binghong(Leo) Li}
\date{\vspace{-0.4in}}

\begin{document}
\maketitle

\section{Problem and Motivation}

LLM-controlled robots selecting actions directly from sensor inputs require memory for extended horizons. Without it, they cannot recall observations, learn from failures, or build persistent world models. Recent LLM memory architectures (MemGPT, Reflexion, Generative Agents) are designed for text-native environments at human speeds. Conversely, robotic systems face fundamentally different constraints: \textbf{Sensor modality} (high-frequency, non-text data streams require continuous processing); \textbf{Strict latency} (memory retrieval sits on the critical path of inference, where delays destabilize control loops); and \textbf{Context scaling} (accumulated long-horizon data must be progressively abstracted and pruned to prevent token overflow). Consequently, naive approaches—appending observations in text or using flat vector retrieval over unprocessed data—fail at scale. A structured approach is necessary but remains empirically unvalidated.

\section{Related Work}

\textbf{LLM agent memory} systems (MemGPT \cite{packer2023memgpt}, Generative Agents \cite{park2023generative}, Reflexion \cite{shinn2023reflexion}) operate purely on text, ignoring sensor translation and real-time latency. \textbf{Embodied memory} integrates semantic forests (Embodied-RAG \cite{xie2024embodiedrag}), dual-tier RAG (PragmaBot \cite{pragmabot2025}), or VLM captioning (STaR \cite{star2025}), while others plan without memory \cite{ahn2022saycan, chen2023typefly}. Critically, these rely on \textit{write-once} memory: observations are stored but never consolidated or pruned. As memory scales, retrieval degrades drastically and token budgets are depleted by stale data. Moreover, hardcoded architectures prevent controlled strategy comparisons. \textbf{Our positioning:} \textit{typemem} decouples memory \emph{policy} from \emph{storage}. This unifies prior approaches under one interface—expressing Embodied-RAG as hierarchical observation and PragmaBot as dual-tier injection—enabling direct comparison against active consolidation strategies to test if consolidation yields measurable scaling benefits.

\section{Approach}

We introduce \textbf{typemem}, an abstract framework decomposing robot memory into three functions over a minimal vector store. \textbf{Observation} runs continuously, converting raw sensor data into natural language memories. \textbf{Consolidation} runs asynchronously, utilizing LLM reasoning to summarize, extract patterns, and prune redundancies. \textbf{Injection} runs synchronously before inference, searching and selecting relevant items within a token budget. By keeping the store unopinionated on retention or linking, all policy lives in these functions, enabling clean ablation. Furthermore, \textbf{LLM-assisted function generation} translates natural language descriptions of a robot's use case into starter Python code and prompts for these functions.

\section{Experimental Plan}

We compare four strategies within \textit{typemem}: \textbf{(1) Full context}, appending raw text up to the token limit with simple FIFO eviction; \textbf{(2) Context compaction}, triggering LLM-based summarization only when out of context space; \textbf{(3) Monolithic RAG}, embedding and storing observations without consolidation for top-$k$ cosine similarity retrieval (KNN); and \textbf{(4) Tiered memory (our proposed approach)}, utilizing active, continuous LLM-generated observation, consolidation, and injection functions. 

\textbf{Evaluation:} \textit{Data generation:} We will collect a continuous offline benchmark by physically patrolling a Unitree Go2 robot dog to gather observation data. \textit{Metrics:} We will evaluate retrieval accuracy and latency (injection time and consolidation overhead) across varying memory scales (100 to 100K entries), measuring performance against manually labeled, ground-truth decision points. \textit{Stretch goal:} Deploying all strategies on the physical quadruped via TypeGo to evaluate real-world task completion and decision quality. 


\newpage
\bibliographystyle{plain}
\begin{thebibliography}{10}

\bibitem{ahn2022saycan}
M.~Ahn, A.~Brohan, N.~Brown, Y.~Chebotar, O.~Cortes, B.~David, C.~Finn, C.~Fu, K.~Gopalakrishnan, K.~Hausman, et~al.
Do as {I} can, not as {I} say: Grounding language in robotic affordances.
In \emph{CoRL}, 2022.
\url{https://arxiv.org/abs/2204.01691}

\bibitem{packer2023memgpt}
C.~Packer, S.~Wooders, K.~Lin, V.~Fang, S.~G. Patil, I.~Stoica, and J.~E. Gonzalez.
{MemGPT}: Towards {LLMs} as operating systems.
\emph{arXiv preprint arXiv:2310.08560}, 2023.
\url{https://arxiv.org/abs/2310.08560}

\bibitem{park2023generative}
J.~S. Park, J.~C. O'Brien, C.~J. Cai, M.~R. Morris, P.~Liang, and M.~S. Bernstein.
Generative agents: Interactive simulacra of human behavior.
In \emph{UIST}, 2023.
\url{https://arxiv.org/abs/2304.03442}

\bibitem{shinn2023reflexion}
N.~Shinn, F.~Cassano, E.~Berman, A.~Gopinath, K.~Narasimhan, and S.~Yao.
Reflexion: Language agents with verbal reinforcement learning.
In \emph{NeurIPS}, 2023.
\url{https://arxiv.org/abs/2303.11366}

\bibitem{chen2023typefly}
G.~Chen, X.~Yu, and L.~Zhong.
{TypeFly}: Flying drones with large language model.
\emph{arXiv preprint arXiv:2312.14950}, 2023.
\url{https://arxiv.org/abs/2312.14950}

\bibitem{star2025}
M.~Yuan, H.~Zhang, M.~Mohammadi, R.~Li, J.~Shan, and S.~L. Waslander.
{STaR}: Scalable task-conditioned retrieval for long-horizon multimodal robot memory.
\emph{arXiv preprint arXiv:2602.09255}, 2026.
\url{https://arxiv.org/abs/2602.09255}

\bibitem{xie2024embodiedrag}
Q.~Xie, S.~Y. Min, P.~Ji, Y.~Yang, T.~Zhang, K.~Xu, A.~Bajaj, R.~Salakhutdinov, M.~Johnson-Roberson, and Y.~Bisk.
{Embodied-RAG}: General non-parametric embodied memory for retrieval and generation.
\emph{arXiv preprint arXiv:2409.18313}, 2024.
\url{https://arxiv.org/abs/2409.18313}

\bibitem{pragmabot2025}
K.~Qu, G.~Lan, R.~Zurbr\"{u}gg, C.~Chen, C.~E. Mower, H.~Bou-Ammar, and M.~Hutter.
A pragmatist robot: Learning to plan tasks by experiencing the real world.
\emph{arXiv preprint arXiv:2507.16713}, 2025.
\url{https://arxiv.org/abs/2507.16713}

\end{thebibliography}

\section*{AI Usage Statement}
AI tools are used to only help the writing of this document. Google Gemini was used to proof read and cut lengths of the proposal.

\end{document}
